\documentclass[10pt,a4paper]{amsart}
\usepackage[margin=19mm]{geometry}
\usepackage{amsmath,amssymb,mathtools}
\usepackage[scr=rsfs,cal=euler]{mathalfa}
\usepackage{xfrac}
\usepackage[unicode,hidelinks,bookmarks=false]{hyperref}
\newcommand{\ie}{i.e.\ }
\newcommand{\cf}{cf.\ }
\newcommand{\resp}{respectively\ }
\newcommand{\Subsection}{\subsection}
\providecommand{\nobreakdash}{\nobreak\hbox{-}}
\setlength{\emergencystretch}{2em}
\multlinegap0pt
\usepackage{amssymb}
\usepackage{xcolor}
\usepackage{tikz}
\usepackage{float}
\usepackage{tcolorbox}
\usepackage{caption}
\newtheorem{lem}{Lemma}
\title{Leaky Zero Forcing on Induced Subgraphs of $d$-dimensional Grid Graphs with an Application to Hopi Rectangles}
\author{Ryan Moruzzi, Jr. \and Sagar Shah \and Aaditeya Tripathi}
\date{Source: arXiv:2509.21529v2 (CC BY 4.0)}
\begin{document}
\maketitle

\noindent\emph{Source and licence.} Leaky Zero Forcing on Induced Subgraphs of $d$-dimensional Grid Graphs with an Application to Hopi Rectangles. Version arXiv:2509.21529v2 (CC BY 4.0), distributed by arXiv under the Creative Commons Attribution 4.0 International licence, http://creativecommons.org/licenses/by/4.0/, as recorded in the arXiv licence metadata for this exact version and shown on the abstract page https://arxiv.org/abs/2509.21529v2. Full licence notice: \texttt{licenses/arxiv-2509.21529-CC-BY-4.0.txt}.\par\smallskip

\noindent\textit{Editorial scope.} Counted unit: the article's lem lem:2d-exposed-neighbors (source0.tex lines 343--356). The article's complete original proof of it is delivered with it (source0.tex lines 358--408, one \texttt{proof} environment of 51 lines). No mathematics is written, repaired or re-derived: the statement and the proof are the article's own lines, in the article's own notation, and no retained line is substituted or re-flowed.  No further source-local context is needed: every cross-reference of every retained line resolves inside the two delivered spans.  Nothing was omitted from any retained span, so the package carries no omission record.  No retained line contains a citation, so the wrapper carries no bibliography and no bibliography entry.  No retained line contains a figure environment, an image-inclusion command or drawing code, so no image is delivered and the package declares no asset dependency.  Editorial formatting only, in addition: an AMS \texttt{amsart} preamble that restates the article's own declarations and macros used by the retained lines, namely the environments \texttt{lem} on separate counters, together with the standard packages those lines need.  The retained environments print in this document as Lemma 1 in order of appearance, whereas the article prints the same items with the numbering of its own full document.  The complete native source of the article as distributed in the arXiv e-print for this exact version is bundled unchanged as \texttt{sources/185751/source0.tex}.
\begin{lem}\label{lem:2d-exposed-neighbors}
Let $d\ge 1$ and let $G = P_{n_1}\square \cdots \square P_{n_d}$. Let $H$ be an induced subgraph of $G$, and set
\[
R \;=\; \{v\in V(H): \deg_H(v)=2d\}.
\]
If $R'\subseteq R$ is a nonempty subset of vertices, then there exist at least $2d$ \emph{distinct} vertices
\[
y_1^+,y_1^-,\dots,y_d^+,y_d^- \in V(H)\setminus R'
\]
such that
\[
\bigl|N_H(y_i^\pm)\cap R'\bigr|=1 \qquad \text{for each } i\in\{1,\dots,d\}.
\]
\end{lem}

\begin{proof}
Assume $R'\subseteq R$ is a nonempty subset of vertices of the set $R \;=\; \{v\in V(H): \deg_H(v)=2d\}$ where $H$ is an induced subgraph of the $d$-dimensional grid graph $G$. Identify the vertices in $G$ with $\{1,\dots,n_1\}\times\cdots\times\{1,\dots,n_d\}\subset \mathbb{Z}^d$ in the usual way, so that
two vertices are adjacent in $G$ if and only if they differ by $\pm \mathbf{e}_i$ in exactly one coordinate, where
$\mathbf{e}_i$ is the $i$th standard basis vector in $d$-dimensional space.

Let $i\in\{1,\dots,d\}$ and choose $x_i^+\in R'$ so that the $i$th coordinate $(x_i^+)_i$ is maximal among vertices in $R'$. Since $\deg_H(x_i^+)=2d$, all neighbors of $x_i^+$ must lie in $H$, and in particular, the neighbor
\[
y_i^+ \;:=\; x_i^+ + \mathbf{e}_i
\]
belongs to $V(H)$. Moreover $y_i^+\notin R'$ by maximality of $(x_i^+)_i$.

We now claim that $N_H(y_i^+)\cap R'=\{x_i^+\}$. Assuming $w\in N_H(y_i^+)\cap R'$, we must have $w=y_i^+\pm \mathbf{e}_j$ for some $j$.
If $j=i$, then $w=y_i^+-\mathbf{e}_i=x_i^+$ or $w=y_i^++\mathbf{e}_i$. By the maximality of the $i$th coordinate $(x_i^+)_i$ in $R'$, we cannot have $w = y_i^++\mathbf{e}_i$, and thus, $w=y_i^+-\mathbf{e}_i=x_i^+$. If $j\neq i$, then $w$ has as its $i$th coordinate $(y_i^+)_i=(x_i^+)_i+1$, again
contradicting maximality of $x_i^+$. Hence, we get $N_H(y_i^+)\cap R'=\{x_i^+\}$.

Similarly, choose $x_i^-\in R'$ so that $(x_i^-)_i$ is minimal among vertices of $R'$. Since $\deg_H(x_i^-)=2d$, the neighbor
\[
y_i^- \;:=\; x_i^- - \mathbf{e}_i
\]
belongs to $V(H)$, satisfies $y_i^-\notin R'$, and an analogous argument gives $N_H(y_i^-)\cap R'=\{x_i^-\}$. Repeating this argument for each $i\in \{1, ... , d\}$ gives $2d$ vertices  $y_1^+,y_1^-,\dots,y_d^+,y_d^-$. 

Lastly, to see that these $2d$ vertices are distinct, fix $i\in\{1,\dots,d\}$ and set
\[
M_i \;:=\; \max\{(v)_i : v\in R'\}
\qquad\text{and}\qquad
m_i \;:=\; \min\{(v)_i : v\in R'\},
\]
where we view $V(G)\subset\mathbb{Z}^d$ as above and $(v)_i$ as the $i$th coordinate of $v$. By construction, $(x_i^+)_i=M_i$ and $(x_i^-)_i=m_i$, so
\[
(y_i^+)_i=(x_i^+)_i+1=M_i+1
\qquad\text{and}\qquad
(y_i^-)_i=(x_i^-)_i-1=m_i-1,
\]
and in particular $y_i^+\neq y_i^-$.

Now let $j\neq i$. Since $y_j^\pm=x_j^\pm\pm \mathbf e_j$ only changes the $j$th coordinate, we have
\[
(y_j^\pm)_i \;=\; (x_j^\pm)_i \text{ for all } i\ne j.
\]
Because $x_j^\pm\in R'$, it follows that $m_i \le (x_j^\pm)_i \le M_i$, and hence
\[
(y_j^\pm)_i \;\le\; M_i \;<\; M_i+1 \;=\; (y_i^+)_i
\qquad\text{and}\qquad
(y_j^\pm)_i \;\ge\; m_i \;>\; m_i-1 \;=\; (y_i^-)_i.
\]
Thus $y_j^\pm \neq y_i^+$ and $y_j^\pm \neq y_i^-$. Since this holds for every $j\neq i$,
the vertices $y_1^+,y_1^-,\dots,y_d^+,y_d^-$ are distinct, and the result follows.


% For instance, if $y_i^+=y_j^+$ with $i\neq j$, then comparing the $i$th coordinate gives $(x_j^+)_i=(x_i^+)_i+1$, contradicting maximality of $(x_i^+)_i$. The other coincidences are ruled out by the
\end{proof}

\end{document}
